%% =====================================================================
\chapter{Appointments}\label{ch:appointments}

\textbf{What it does.} Patients see a doctor's free slots and book one. Doctors
confirm, complete or mark no-show. Patients can reschedule; either side can
cancel. Completing a visit raises the bill.

\textbf{Files.} \file{src\codesl{}models\codesl{}Appointment\codedot{}js},
\file{src\codesl{}modules\codesl{}appointments\codesl{}slots\codedot{}js},
\file{src\codesl{}modules\codesl{}appointments\codesl{}appointments\codedot{}routes\codedot{}js},
\file{src\codesl{}modules\codesl{}users\codesl{}users\codedot{}routes\codedot{}js} (available slots).

\section{Generating slots from availability}

Doctors do not create slots one by one. Each doctor has an availability rule
(for example Monday--Friday, 9:00--17:00, 30-minute slots) and slots are
computed from it.

\begin{lstlisting}[language=JavaScript,caption={buildDaySlots and isValidSlot (\file{src/modules/appointments/slots.js})}]
function buildDaySlots(availability, dateString, maxAdvanceDays, now = new Date()) {
  const [year, month, day] = dateString.split('-').map(Number);
  const date = new Date(year, month - 1, day);
  if (Number.isNaN(date.getTime())) return [];

  const latest = new Date(now);
  latest.setDate(latest.getDate() + maxAdvanceDays);
  if (date > latest) return [];
  if (!availability.days.includes(date.getDay())) return [];

  const slots = [];
  const stepMs = availability.slotMinutes * 60 * 1000;
  const dayEnd = new Date(year, month - 1, day, availability.endHour, 0, 0, 0);
  for (let t = new Date(year, month - 1, day, availability.startHour, 0, 0, 0); t.getTime() + stepMs <= dayEnd.getTime(); t = new Date(t.getTime() + stepMs)) {
    if (t > now) slots.push({ start: t, end: new Date(t.getTime() + stepMs) });
  }
  return slots;
}

/** True when `start` is exactly one of the doctor's slot boundaries. */
function isValidSlot(availability, start, maxAdvanceDays, now = new Date()) {
  const dateString = `${start.getFullYear()}-${String(start.getMonth() + 1).padStart(2, '0')}-${String(start.getDate()).padStart(2, '0')}`;
  return buildDaySlots(availability, dateString, maxAdvanceDays, now).some((s) => s.start.getTime() === start.getTime());
}
\end{lstlisting}

\begin{howto}
\begin{itemize}
  \item Days outside the doctor's working days, or beyond the hospital's
    booking window (\code{maxAdvanceBookingDays}, an admin setting), have no
    slots.
  \item A slot is included only if it ends by \code{endHour}, and only if it is
    still in the future.
  \item \code{now} is a parameter so the unit tests can fix the clock.
  \item \code{isValidSlot} is what the booking endpoint uses: a client cannot
    book 9:07 or 3 a.m. by crafting a request, only a real slot.
\end{itemize}
\end{howto}

\code{GET \codesl{}api\codesl{}doctors\codesl{}:doctorId\codesl{}available-slots?date=YYYY-MM-DD} builds the
day's slots, then looks up appointments that already hold one of them, and marks
each slot \code{available: true\codesl{}false}. The patient's booking dialog shows taken
slots greyed out.

\section{Double-booking prevention}

Hiding taken slots in the interface is not enough: two patients can open the
dialog at the same moment and click the same slot. The guarantee is enforced by
the database:

\begin{lstlisting}[language=JavaScript,caption={Partial unique index (\file{src/models/Appointment.js})}]
  // True while the appointment occupies its slot. Cancelling sets it to false,
  // which frees the slot for someone else.
  holdsSlot: { type: Boolean, default: true },
}, { timestamps: true });

// Double-booking prevention at the database level: among appointments that
// hold their slot, a doctor can have only one per start time. Even two
// requests racing each other cannot both succeed.
appointmentSchema.index(
  { doctorId: 1, startTime: 1 },
  { unique: true, partialFilterExpression: { holdsSlot: true } }
);

appointmentSchema.pre('validate', function syncHoldsSlot(next) {
  this.holdsSlot = ACTIVE_STATUSES.includes(this.status);
  next();
});
\end{lstlisting}

\begin{howto}
A \emph{unique index} makes MongoDB reject a second document with the same
\code{(doctorId, startTime)}. A \emph{partial} index applies only to documents
matching a filter: here, those with \code{holdsSlot: true}. Pending, confirmed
and completed appointments hold their slot; cancelled and no-show ones do not.
So cancelling an appointment frees the slot automatically, while a race between
two bookings always ends with exactly one success. The \code{pre('validate')}
hook keeps \code{holdsSlot} in step with \code{status} on every save.
\end{howto}

\section{Booking: \code{POST \codesl{}api\codesl{}appointments}}

\begin{lstlisting}[language=JavaScript,caption={Booking a slot (\file{src/modules/appointments/appointments.routes.js})}]
router.post('/', requireRole('patient'),
  validate({ body: Joi.object({
    doctorId: objectId().required(),
    startTime: Joi.date().iso().greater('now').required(),
    type: Joi.string().valid('consultation', 'follow_up', 'checkup').default('consultation'),
    reason: Joi.string().trim().min(3).max(500).required(),
  }) }),
  asyncHandler(async (req, res) => {
    const { doctorId, startTime, type, reason } = req.body;
    const doctor = await User.findOne({ _id: doctorId, role: 'doctor', approvalStatus: 'approved', isActive: true });
    if (!doctor) throw notFound('Doctor not found');

    const settings = await HospitalSettings.getSingleton();
    const start = new Date(startTime);
    if (!isValidSlot(doctor.doctorProfile.availability, start, settings.maxAdvanceBookingDays)) {
      throw badRequest('That time is not one of the doctor\'s available slots');
    }

    let appointment;
    try {
      appointment = await Appointment.create({
        patientId: req.user.userId,
        doctorId,
        startTime: start,
        endTime: new Date(start.getTime() + doctor.doctorProfile.availability.slotMinutes * 60000),
        type,
        reason,
      });
    } catch (error) {
      if (error.code === 11000) throw conflict('This slot has just been booked. Please choose another time.', 'SLOT_TAKEN');
      throw error;
    }

    await recordAppointmentEvent(appointment, 'CREATE', req);
    await notificationService.notify(doctorId, {
      type: 'info', title: 'New appointment request',
      message: `A patient requested a ${type.replace('_', '-')} on ${fmt(start)}: "${reason}"`, page: 'appointments',
    });
    // ... 201 with the populated appointment
\end{lstlisting}

\begin{howto}
Only verified, active doctors can be booked. The duplicate-key error (11000)
from the unique index is translated into a friendly 409 \code{SLOT\_TAKEN}; the
frontend shows the message and reloads the slot grid. Every booking is a ledger
entry and a notification for the doctor.
\end{howto}

\section{Status changes, rescheduling and cancelling}

\begin{figure}[H]
\centering
\begin{tikzpicture}[node distance=1.0cm and 1.6cm]
  \node[box] (p) {pending};
  \node[box,right=of p] (c) {confirmed};
  \node[box,right=of c,draw=okgreen,fill=green!6] (d) {completed\\\scriptsize raises a bill};
  \node[box,below=of c] (n) {no\_show};
  \node[box,below=of p,draw=badred,fill=red!3] (x) {cancelled};
  \draw[arr] (p) -- node[lbl,above]{doctor} (c);
  \draw[arr] (c) -- node[lbl,above]{doctor} (d);
  \draw[arr] (c) -- node[lbl,right]{doctor} (n);
  \draw[arr] (p) -- node[lbl,left]{either} (x);
  \draw[arr] (c) -- node[lbl,above,sloped]{either} (x);
  \draw[arr] (c.north) to[bend right=35] node[lbl,above]{patient reschedules} (p.north);
\end{tikzpicture}
\caption{Appointment states}\label{fig:appt-states}
\end{figure}

\begin{lstlisting}[language=JavaScript,caption={Allowed doctor transitions (\file{appointments.routes.js})}]
const DOCTOR_TRANSITIONS = { pending: ['confirmed'], confirmed: ['completed', 'no_show'] };
// ...
      const allowed = DOCTOR_TRANSITIONS[appointment.status] || [];
      if (!allowed.includes(req.body.status)) {
        throw conflict(`Cannot change an appointment from "${appointment.status}" to "${req.body.status}"`, 'INVALID_TRANSITION');
      }
      if (req.body.status === 'completed' && appointment.startTime > new Date()) {
        throw conflict('An appointment cannot be completed before it starts', 'INVALID_TRANSITION');
      }
      appointment.status = req.body.status;
      if (req.body.notes !== undefined) appointment.notes = req.body.notes;
      await appointment.save();
      await recordAppointmentEvent(appointment, req.body.status.toUpperCase(), req);
      // ... notify the patient
      if (req.body.status === 'completed') {
        const doctor = await User.findById(appointment.doctorId);
        await createBillForAppointment(appointment, doctor, req.user);
      }
\end{lstlisting}

\begin{itemize}
  \item \textbf{Who may do what} is checked on the server: only the doctor
    changes status (and only once approved), only the patient reschedules.
  \item \textbf{Rescheduling} (\code{PATCH} with \code{startTime}) re-runs
    \code{isValidSlot}, moves the appointment, sets it back to \emph{pending}
    (the doctor must confirm the new time) and catches the same 11000 error if
    the new slot was just taken.
  \item \textbf{Cancelling} (\code{DELETE}) is allowed for pending and confirmed
    appointments; the record stays (with who cancelled and why) and the slot is
    freed through \code{holdsSlot}.
  \item An appointment that is not yours returns \textbf{404}, not 403, so
    appointment ids cannot be probed.
\end{itemize}
